% Chapter 15: BSP
% File: chapters/chapter-15-bsp/chapter-15-bsp.tex

\chapter{Board Support Packages and Hardware Abstraction}
\label{ch:chapter-15}
\chapthumb{15}

\begin{quote}
    \textit{``A BSP is the layer that knows which wire is which — and it is the only layer allowed
    to.''}
\end{quote}

The PAC knows registers. The HAL knows protocols. Neither knows that a particular board puts an LED on
\codeinline{PA7} with the cathode to the pin, and that is the BSP's job: the last layer between the
user's intent and the silicon, and the only one that should contain a pin number.

This chapter builds a BSP for hardware that arrives in two pieces — a module and a carrier — which is
the arrangement that exposes every mistake in the design.

\minitoc

\clearpage

\section{Two Pieces of Hardware, Not One}
\label{sec:bsp-two-pieces}

This board is not one board. It is a module — the XIAO MG24 Sense, which runs on its own — plus a
carrier that the module plugs into. The distinction is physical, it decides who owns which feature, and
a single BSP crate erases it.

\subsection{What Each Piece Owns}

\begin{table}[htbp]
    \centering
    \footnotesize
    \begin{tabular}{@{}p{0.22\textwidth}p{0.30\textwidth}p{0.40\textwidth}@{}}
        \toprule
        \textbf{Piece} & \textbf{Components} & \textbf{Can be used alone?} \\
        \midrule
        XIAO MG24 & LED \codeinline{PA7}, USB, radio & yes \\
        XIAO MG24 \emph{Sense} & the above + IMU \codeinline{LSM6DS3TR-C}, analog mic & yes \\
        Expansion Base & OLED, BME68x, RTC, buzzer, button, Grove, microSD & no; needs a XIAO \\
        \bottomrule
    \end{tabular}
    \caption{Three hardware identities and one board on the desk. The Sense is a variant of the module,
    not a different module.}
    \label{tab:bsp-pieces}
\end{table}

\subsection{The Same PCB With More Parts}

The Sense is a XIAO with two extra components fitted. That single fact decides the crate layout:

\begin{itemize}
    \item \textbf{A feature, not a crate.} Two crates would each need the boot sequence, the clock
          tree, the LED and the pin map — four things to keep in step, for two products that are
          physically the same PCB.
    \item \textbf{A feature that is opt-in.} A user with a plain XIAO must not compile the IMU driver
          or its bus. \codeinline{--features sense} adds exactly those.
    \item \textbf{The feature names the thing that is fitted,} not the price bracket or the marketing
          name: \codeinline{sense} means ``the two sensors are populated''.
\end{itemize}

\begin{minted}{toml}
[features]
default = []
# The Sense variant is the same PCB with more components fitted: the IMU and the
# analog microphone. Same board, fewer parts -- so a feature, not a crate.
sense = []
\end{minted}

\subsection{The Carrier Is a Different Crate, and It Points One Way}

The Expansion Base adds devices the module has never heard of, and a user may own neither, one or both.
So it is its own crate, and the dependency has a direction:

\begin{minted}{text}
xiao-expansion-base-bsp  ──depends on──▶  xiao-mg24-bsp  ──depends on──▶  efr32mg24-hal
\end{minted}

The opposite arrow — the module crate knowing that a carrier exists — is the defect this book's
section~\ref{sec:layer-boundaries} describes: it makes the module unusable without the carrier's code,
and it puts carrier pins in a crate whose whole claim is ``the XIAO on its own''.

\subsection{What the Direction Buys}

\begin{itemize}
    \item A user with only the XIAO compiles one small crate, not the carrier's drivers.
    \item Carrier examples live where the carrier is, so the path of a file tells you the hardware it
          needs before you open it.
    \item The module crate cannot drift into knowing the carrier, because the manifest forbids it —
          and \codeinline{cargo} enforces manifests even when reviewers do not.
\end{itemize}

\section{One Crate per Board, and the Arrow Between Them}
\label{sec:bsp-one-crate}

The crate layout is not a matter of taste. It follows from one question: \emph{which of these things
can be owned separately?}

\subsection{The Two Schools}

\begin{table}[htbp]
    \centering
    \footnotesize
    \begin{tabular}{@{}p{0.24\textwidth}p{0.32\textwidth}p{0.36\textwidth}@{}}
        \toprule
        \textbf{Project} & \textbf{Layout} & \textbf{Where the wiring lives} \\
        \midrule
        \codeinline{stm32f4xx-hal} & one HAL per family, no BSP & in every example \\
        \codeinline{esp-hal} & one HAL, no BSP, chip by feature & in every example \\
        \codeinline{rp-hal} / \codeinline{rp-hal-boards} & HAL per chip + a BSP crate per board & in the BSP crate \\
        \bottomrule
    \end{tabular}
    \caption{Two schools. The difference is not sophistication: it is whether the board's wiring is a
    fact worth writing down once.}
    \label{tab:bsp-schools}
\end{table}

Both schools are legitimate. The first optimises for a HAL whose users all have a schematic in front of
them; the second optimises for a board whose wiring was expensive to establish. This project is firmly
in the second camp, for a reason that is in chapter~4's table of contradictions: the pin map of this
board was measured, and a measurement that is repeated in every example is a measurement that will be
repeated wrongly.

\subsection{The Four Crate Questions}

\begin{enumerate}
    \item \textbf{How many chips?} One HAL per chip or per family. This part gets one.
    \item \textbf{Do the boards stack?} If a carrier holds a module, one crate per piece, arrow carrier
          $\rightarrow$ module.
    \item \textbf{Are there variants?} A variant of the \emph{same} PCB is a feature. A different PCB is
          a crate.
    \item \textbf{Was the wiring obvious or was it conquered?} Obvious wiring can live in examples;
          conquered wiring belongs in a BSP with the measurement next to it.
\end{enumerate}

\subsection{What Ends Up in the Workspace}

\begin{minted}{text}
crates/
  efr32mg24-pac/            registers
  efr32mg24-hal/            drivers + pin and bus primitives
  xiao-mg24-bsp/            Module: led(), and with --features sense: imu(), microphone()
  xiao-expansion-base-bsp/  ExpansionBase: oled(), environment(), rtc(), buzzer(), button()
libraries/
  motion-detection/         domain policy: no hardware at all
\end{minted}

Two rules hold that shape together, and both are checks in this project's gate:

\begin{itemize}
    \item \textbf{Nothing in libraries/ may depend on the hardware stack.} A classifier
          that needs a HAL is a classifier that cannot be tested on a laptop, and the test suite is the
          only thing that makes it trustworthy.
    \item \textbf{Every crate that exposes hardware carries at least one example of its own level.} An
          API with no examples compiles, lints, and rots — this project had a whole module do exactly
          that for weeks.
\end{itemize}

\subsection{The Defect the Layout Prevented}

Before the split, the HAL itself contained \codeinline{board.rs} with the carrier's wiring — the OLED
bus, the button pin, the buzzer channel — because it was convenient. Two consequences followed
immediately:

\begin{enumerate}
    \item The HAL became unusable without the carrier's assumptions, so a user with a bare XIAO had to
          read carrier pins to blink an LED.
    \item Two BSPs existed at once: the HAL's board module and the new facade. Half the examples proved
          one, half the other, and the older one stopped compiling without anyone noticing.
\end{enumerate}

The fix was not better documentation. It was deleting the wiring from the HAL and making the
dependency direction a manifest fact.

\section{What a BSP Should Expose}
\label{sec:bsp-api}

A BSP's public API is a vocabulary, and the test of a vocabulary is what a user can say with it without
knowing anything else.

\subsection{The Right Shape}

\begin{minted}{rust}
let module = Module::take().expect("peripherals taken twice");
let mut led = module.led();

let base = ExpansionBase::take().expect("peripherals taken twice");
let mut oled = base.oled();
let env = base.environment()?;      // BME68x
let rtc = base.rtc()?;              // PCF8563
let mut buzzer = base.buzzer();
\end{minted}

Every one of those names is a thing on the desk. None of them is an address, a bus, a pin, an alias or
a register — and that is the whole specification.

\subsection{What the Signature Hides, and What It Must Not}

\begin{table}[htbp]
    \centering
    \footnotesize
    \begin{tabular}{@{}p{0.26\textwidth}p{0.30\textwidth}p{0.40\textwidth}@{}}
        \toprule
        \textbf{Hidden} & \textbf{Example} & \textbf{Why hiding it is right} \\
        \midrule
        addresses & \codeinline{0x3C}, \codeinline{0x76}, \codeinline{0x51} & a user cannot check them and must not need to \\
        pins & \codeinline{PC04}/\codeinline{PC05}, \codeinline{PB2}/\codeinline{PB3} & measured facts; wrong in half the documents \\
        power gates & \codeinline{PD5}, \codeinline{PC8} & sequencing the user would get wrong \\
        aliases & which peripheral window answered & a security-attribution detail \\
        baud and clock & I2C half-period, PWM frequency & driver knowledge \\
        \bottomrule
    \end{tabular}
    \caption{What a facade should hide. Note that every row is a \emph{fact}, not a \emph{choice} — the
    user has no decision to make about any of them.}
    \label{tab:bsp-hides}
\end{table}

\subsection{What It Must Not Hide: Structure}

The temptation is to hide the two buses, because ``two I2C buses'' is an implementation detail. It is
not — it is a fact about the hardware that changes what the user can do:

\begin{itemize}
    \item The panel, the BME68x and the RTC share \codeinline{PC04}/\codeinline{PC05}.
    \item The IMU has its own pair, \codeinline{PB2}/\codeinline{PB3}, because the module's sensors are
          wired to the module's pins.
    \item A single \codeinline{i2c()} would have to take a ``which bus'' argument anyway, so hiding it
          only moves the choice to a place with less information.
\end{itemize}

The rule that came out of this: \textbf{hide the addresses, keep the structure}. A facade that hides
structure has to guess, and guessing is what the layer below exists to prevent.

\subsection{Conflicts Should Be Refused, Not Resolved}

\begin{minted}{rust}
// The microSD card's chip select shares a pin on this board. When that driver lands,
// this facade is where the conflict is refused -- not silently arbitrated.
pub fn sd(&mut self) -> Result<SdCard<'_>, BoardError>
\end{minted}

A facade that resolves a real conflict silently is a facade that produces a bug the user cannot see.
Returning an error — or not offering the operation at all — is honest, and it is the difference between
an API that documents the hardware and one that pretends the hardware is simpler than it is.

\section{Examples Belong to the Hardware They Need}
\label{sec:bsp-examples}

The most practical BSP rule in this book is about where example files live, because that is where a
layout either documents itself or lies.

\subsection{The Rule}

\begin{quote}
    An example lives in the crate whose \textbf{smallest} hardware requirement it satisfies.
\end{quote}

\begin{table}[htbp]
    \centering
    \footnotesize
    \begin{tabular}{@{}p{0.36\textwidth}p{0.28\textwidth}p{0.32\textwidth}@{}}
        \toprule
        \textbf{Example} & \textbf{Lives in} & \textbf{Needs} \\
        \midrule
        \codeinline{module\_led\_blink} & \codeinline{xiao-mg24-bsp} & the module \\
        \codeinline{module\_imu} & \codeinline{xiao-mg24-bsp} (\codeinline{sense}) & the module, Sense variant \\
        \codeinline{base\_env\_clock} & \codeinline{xiao-expansion-base-bsp} & module + carrier \\
        \codeinline{base\_fall\_detector} & \codeinline{xiao-expansion-base-bsp} & module + carrier (panel, buzzer) \\
        \bottomrule
    \end{tabular}
    \caption{The path of the file is the bill of materials.}
    \label{tab:bsp-examples}
\end{table}

\subsection{The Three Ways to Get It Wrong}

\begin{enumerate}
    \item \textbf{Everything in the carrier crate.} Convenient, and it makes the module crate's API
          unexercised: this project's module facade had no examples for weeks, and its documentation
          described a pin that did not exist.
    \item \textbf{Everything in the module crate.} Impossible without the carrier's devices, so the
          examples would need \codeinline{cfg} guards and would prove neither level.
    \item \textbf{A flat examples/ directory at the workspace root.} The path stops telling you
          anything, and \codeinline{cargo} stops checking that the crate still builds its own examples.
\end{enumerate}

\subsection{The Check That Keeps It True}

\begin{minted}{bash}
# every crate that exposes hardware must carry at least one example of its own level
for dir in crates/*/ libraries/*/; do
  [ -f "${dir}Cargo.toml" ] || continue
  n=$(find "$dir/examples" -maxdepth 1 -name '*.rs' | wc -l)
  [ "$n" -gt 0 ] || fail "$(basename "$dir"): no examples -- its API is not built by anything"
done
\end{minted}

That check is why this project's module facade has a \codeinline{module\_led\_blink} that uses only the
module, and why the carrier's examples are in the carrier. A rule enforced by the file system is a rule
nobody has to remember.

\subsection{The Half-Example That Taught the Rule}

A fall detector was written in the HAL's examples directory. It used the IMU, the panel and the buzzer:
two boards' worth of hardware, in the crate that knows neither. It compiled, because the HAL happened to
contain the wiring at the time — which was itself the defect the split fixed.

After the split, the same program had an obvious home: the carrier crate, which depends on the module
and therefore has both the IMU (through the module) and the panel (its own). Moving it changed no code.
That is the sign of a layout that was right: the program did not have to change to be correctly placed.


\clearpage

\begin{keytakeaways}
    \item Hardware that can be owned in pieces needs a crate per piece, with the dependency arrow
          pointing one way: carrier depends on module, module depends on HAL.
    \item A variant of the same PCB is a feature, not a crate. Two crates would duplicate the boot,
          the clock tree and the LED, and one of them would drift.
    \item A BSP's API is what a user asks for — \codeinline{led()}, \codeinline{imu()},
          \codeinline{oled()} — never an address, an alias or a register.
    \item A facade that hides a real conflict (two buses, a shared pin) hides it until the user hits
          it. Hide the addresses; keep the structure visible.
    \item Every example belongs to the crate whose hardware it needs, and the gate fails a crate that
          has hardware and no example.
\end{keytakeaways}
